\documentclass[11pt,a4paper]{article}
\usepackage[T1]{fontenc}
\usepackage{lmodern,amsmath,amssymb,amsthm,mathtools,microtype}
\usepackage[margin=27mm]{geometry}
\usepackage[colorlinks=true,linkcolor=blue,citecolor=blue,urlcolor=blue]{hyperref}
\setlength{\emergencystretch}{2em}
\newtheorem{theorem}{Theorem}[section]
\newtheorem{lemma}[theorem]{Lemma}
\newtheorem{proposition}[theorem]{Proposition}
\newtheorem{corollary}[theorem]{Corollary}
\newcommand{\Q}{\mathbb Q}\newcommand{\Z}{\mathbb Z}
\DeclareMathOperator{\Tr}{Tr}\DeclareMathOperator{\N}{N}
\title{A scalar normalization for local square-complement realization}
\author{}\date{16 November 2026}
\begin{document}\maketitle
\begin{abstract}
We give a scalar coefficient for the general elliptic detector source
whose complement is identically a nonzero rational square. The source
retains its twenty-five-coordinate presentation and its soundness for
every sufficiently large odd coefficient exponent, under the same
arithmetic premises. A depth--height inequality controls the fixed
coefficient exponent. For every integer input, the actual guarded
detector tuples can then have square complements at any prescribed finite
set of rational primes. This is a finite-local realization theorem; it
does not establish rational square complements, a common rational pole
fiber, a smaller defining polynomial, or an effective sufficient exponent.
\end{abstract}

\section{The modified scalar and its scope}

Put $\ell=76249$. For an odd integer $L\ge3$ and a rational input $X$, set
\begin{align}
 \beta&=216X^4+450459,\label{eq:beta}\\
 r&=\beta(\beta-1)(\beta-2401)(32\beta+343)(2\beta-1),\label{eq:r}\\
 \rho_*&=\frac{102^2\ell^L r}{(1+153\ell^L r)^2}.\label{eq:rho}
\end{align}
Use the general native source of \cite[Sections 3--6]{UnitNorm}, with
$\rho_*$ in place of its scalar coefficient $\rho$ everywhere: in the
comparison scale, the detector weights, and the detector inputs.
All point roles, contacts, difference formulas, finite point blocks,
and pole equations are retained. This is a change of the early scalar
coefficient, before any comparison slope or auxiliary point is chosen.
The separate substitution $\Phi(b)=1156b/(17+b)^2$ on an already
formed pole input does not preserve its denominator obstruction
\cite{Normalization}; the coefficient construction here retains the
original detector products and their local extraction arguments.

In particular the two actual detector inputs remain
\begin{equation}\label{eq:detector}
 b_j=(\rho_*R_j)
       \left(\frac{\N_{K/\Q}(W_j)}{(W_j)_0}\right)^2,
 \qquad K=\Q(\sqrt2),\quad j=0,1,
\end{equation}
where $R_j$ and $W_j$ are the specified native functions with the new
coefficient. Each $R_j$ has the same ten simple contacts as in the source.
The comparison scales are $1$ and $1+\rho_*T^2$, where $T$ is the
rational trace of the sum of the four fixed root abscissas.

The arithmetic premises are exactly those retained in
\cite[Section 1]{General} and \cite[Section 1]{UnitNorm}: the elliptic
group and height assertions, the normalized contact-height realization,
the prime-pattern theorem with its prescribed finite conditions, and the
pole obstruction and prepared-input realization. The aggregate
trace-height bound is used with the same scope as in those sources.

\begin{samepage}
\begin{theorem}\label{thm:main}
Under these retained arithmetic premises, the modified source has the
following properties.
\begin{enumerate}
\item For every integer input and every odd $L\ge3$, it has rational
      witnesses with prepared detector inputs and all native guards
      nonzero at every geometric character of the point algebra.
\item For every sufficiently large fixed odd $L$, its rational input
      set is exactly $\Z$.
\item For any prescribed finite set $S$ of rational primes, the witnesses
      in part~1 can additionally satisfy
      \[
      17(17-b_j)\in\Q_p^{\times2}\qquad(p\in S,\ j=0,1).
      \]
      Both real complements are positive. There are independent
      auxiliary-index return families preserving all these finite-local
      conditions and all point-stage guards.
\end{enumerate}
The finite-algebra presentation still has twenty-five total rational
coordinates, including the input and its single separating coordinate.
No effective sufficient exponent or improved coefficient bound is asserted.
\end{theorem}
\end{samepage}

The conclusion in part~3 concerns the actual source values in
\eqref{eq:detector}. It does not permit assigning arbitrary prepared
rational inputs to the pole predicates. Neither part~3 nor the independent
return statement imposes a global rational-square equation.

\section{The scalar square identity and local valuations}

\begin{lemma}\label{lem:square}
For every rational $X$ and odd $L\ge3$,
\[
 0<\rho_*<1/17,\qquad
 17(17-\rho_*)=
 \left(\frac{17(1-153\ell^Lr)}{1+153\ell^Lr}\right)^2\ne0.
\]
For integer $X$, the scalar $\rho_*$ is prepared for the original pole
predicate.
\end{lemma}
\begin{proof}
Write $z=153\ell^Lr$. Then $102^2=68\cdot153$ gives
$\rho_*=68z/(1+z)^2$ and the displayed square identity. Since
$\beta\ge450459$ at every real input, all five factors in~\eqref{eq:r}
exceed one. Thus $z>1156$, so $0<\rho_*<68/z<1/17$, and the square
root displayed above is nonzero.

For integer $X$, $r$ is a positive integer divisible by four.
The denominator $1+153\ell^Lr$ is a unit at $2,3,17$.
The original source denominator
$1+17(\ell^{(L+1)/2}+1)r$ is a unit at the same primes, and the ratio
of the two scalar coefficients is the square of the ratio of their
denominators. Therefore the local square classes and valuations at
$2,3,17$ agree with those of the prepared scalar in \cite{General}.

The new numerator and denominator are coprime. A prime dividing
$1+153\ell^Lr$ divides neither $153\ell^L$ nor $r$; it cannot divide
the remaining factor $68$, because the denominator is odd and $17$
divides $153$. Hence the reduced denominator of $\rho_*$ is a square.
This proves its preparation.
\end{proof}

Let $\mathcal S$ be the fixed exceptional set of the general source.
In particular it contains $2,3,17$ and all primes at most $76176$,
but not $\ell$.

\begin{lemma}\label{lem:odd}
Every prime contributing to the contact radical $M_5(\beta)$ outside
$\mathcal S$ has positive odd valuation in $\rho_*$. The scalar-odd
and mixed-comparison conclusions of \cite[Section 3]{General}, with
the unsquared-scale modification of \cite{UnitNorm}, continue to hold.
Consequently, for a source solution with scalar root label $a$ and
$c=|a-1|>0$, there is a fixed $H_1$ such that
\begin{equation}\label{eq:scalar-height}
                         h(\beta)\le H_1c^{7/9}.
\end{equation}
\end{lemma}
\begin{proof}
Fix $p\notin\mathcal S\cup\{\ell\}$. If $v_p(r)>0$, the new
denominator is a unit, so $v_p(\rho_*)=v_p(r)$. If $v_p(r)<0$, then
$v_p(\beta)<0$; it is a multiple of four, and all five raw factors
have that valuation. Thus $v_p(r)=5v_p(\beta)$ is even, and
$v_p(\rho_*)=-v_p(r)$ is even. If $v_p(r)=0$, the valuation of
$\rho_*$ is nonpositive and even. At a finite simple contact of
positive odd order, precisely the first case applies. The fixed residue
certificate at $\ell$ excludes every finite contact of $\beta$ there;
its possible negative valuation is again a multiple of four.

The scalar-odd proof uses the odd valuation of the scalar, the unit
alternative for one contact function, and the detector trace inequality.
All remain valid for $\rho_*$. In particular an odd valuation
$\mu=v_p(\rho_*)$ in a source solution is positive and gives
$v_p(c)\ge2\mu+1\ge3$. The scale $1+\rho_*T^2$ has the same
unit alternative because $v_p(\rho_*T^2)$ is odd. The mixed-comparison
proof already includes negative even scalar valuations. Thus its two
residual divisibilities are unchanged. Applying the same contact-height
theorem to the scalar radical gives~\eqref{eq:scalar-height}.
\end{proof}

\section{A depth--height estimate}

\begin{lemma}\label{lem:tradeoff}
For every rational source solution,
\begin{equation}\label{eq:tradeoff}
             L\log\ell\le\tfrac12\log c+5h(\beta).
\end{equation}
There is a constant $C$, independent of $X,L$, such that
\begin{equation}\label{eq:rho-height}
                 h(\rho_*)\le20h(\beta)+\log c+C.
\end{equation}
Moreover $c$ tends uniformly beyond every fixed bound as $L$ tends to
infinity through odd exponents among source solutions.
\end{lemma}
\begin{proof}
The fixed-prime residue certificate states that $r$ is an $\ell$-unit
when $X$ is $\ell$-integral. If $m=-v_\ell(X)>0$, then
$v_\ell(\beta)=-4m$ and $v_\ell(r)=-20m$. Because $L$ is odd,
the valuation of $153\ell^Lr$ is nonzero, and
\[
 \mu=v_\ell(\rho_*)=
 \begin{cases}
 L,&v_\ell(X)\ge0,\\
 |L-20m|,&m=-v_\ell(X)>0.
 \end{cases}
\]
Thus $\mu$ is positive and odd. Lemma~\ref{lem:odd} gives
$v_\ell(c)\ge2\mu+1$, hence $\mu\log\ell\le\frac12\log c$.
In the second case, $20m\log\ell\le5h(\beta)$. Now
$L\le|L-20m|+20m$ proves~\eqref{eq:tradeoff}; the integral case
is immediate.

For the height estimate, write $r=u/v$ in lowest terms and set
$M=\max(|u|,v)$, $B=153\ell^L$. The expression
\[
                 \rho_*=\frac{68Buv}{(v+Bu)^2}
\]
has numerator and denominator bounded in absolute value by
$(B+1)^2M^2$, since $B\ge66$ and $(B+1)^2\ge68B$.
Reduction of the fraction can only decrease this bound. Hence
\[
 h(\rho_*)\le2h(r)+2L\log\ell+2\log154
              \le10h(\beta)+2L\log\ell+C.
\]
Use~\eqref{eq:tradeoff} to obtain~\eqref{eq:rho-height}.
Finally \eqref{eq:scalar-height} and \eqref{eq:tradeoff} imply
\[
 L\log\ell\le\tfrac12\log c+5H_1c^{7/9}.
\]
The right side is bounded on every fixed finite range of positive
integers $c$, proving the last assertion.
\end{proof}

The stronger estimate $v_\ell(\rho_*)\ge L$ is false. For example
$L=21$ and $X=1/\ell$ give $v_\ell(\rho_*)=1$. The height term in
\eqref{eq:tradeoff} is essential.

\section{Soundness and actual finite-local completeness}

\begin{proof}[Proof of Theorem~\ref{thm:main}]
For soundness, retain the notation of \cite[Sections 3--4]{General}:
$B_{\mathrm{ind}}$ is the largest absolute root label and $N$ is either
the nonzero Lucas residual or the nonzero direct-input residual.
Lemma~\ref{lem:odd} supplies exactly the same two radical divisibilities.
The contact-height and aggregate trace-height arguments, with
\eqref{eq:rho-height}, therefore give fixed constants $C_2,J_2$ such that
\[
 B_{\mathrm{ind}}^2\le
 C_2|Nc|^\eta+C_2c^{7/9}+C_2\log c+J_2,
\]
where $\eta=7/18$ for the Lucas residual and $\eta=7/6$ for the
direct-input residual. The constants are independent of $X,L$.
All exponents in the two residual eliminations are unchanged: the first
would bound $B_{\mathrm{ind}}^2$ by terms of smaller degree, and after
its vanishing the Lucas growth estimate contradicts the second unless
that residual also vanishes. The required lower bounds on $c$ are fixed.
Lemma~\ref{lem:tradeoff} supplies them for every sufficiently large odd
$L$. Thus $\beta$ is an integer; by~\eqref{eq:beta}, integrality of
$216X^4+450459$ for rational $X$ implies $X\in\Z$.

For completeness, fix an integer $X$, odd $L\ge3$, and the prescribed
finite set $S$ (which may be empty).
Lemma~\ref{lem:square} supplies the same preparation and small-real
properties of the scalar required in \cite[Section 5]{General}.
Its support is finite once $X,L$ are fixed. At this stage, before the
scalar functions or root points are chosen, enlarge the preparation set
to contain $S$. Lemma~\ref{lem:square} supplies the nonzero rational
square $C_0=17(17-\rho_*)$. Choose the finite precisions so that
closeness of $b_j$ to $\rho_*$ will make $17(17-b_j)$ a nonzero
square at every prime in $S$. This is possible because
$\Q_p^{\times2}$ is open in $\Q_p^\times$, including at $2$ and $17$.
Now prepare scalar functions,
choose the root labels and comparison scales, select fresh contact
primes and the slope, and finally choose the auxiliary index, in that
order. The inherited prime-pattern theorem permits the fixed support
and any additional finite preparation set to be avoided by the fresh
primes. At those primes, the original norm/trace argument makes odd
valuations of the actual $b_j$ positive. At the scalar support and at
$2,3,17$, use the unchanged closeness argument. The remaining valuations
are even. The real construction gives
$0<b_j<8\rho_*<17$. Thus both actual inputs are prepared.

The same construction makes $R_j,W_{j,+},W_{j,-}$ as close to one
as specified at each $p\in S$. Formula~\eqref{eq:detector} makes the
actual $b_j$ sufficiently close to $\rho_*$, giving the prescribed
local-square complements. Their real positivity follows from $b_j<17$.

The full-character guard proof of \cite[Proposition 5.3]{UnitNorm}
retains all already prescribed finite-place conditions while deleting
the finitely many real guard zero sets. It uses positivity of the scalar,
fixed root data before slope selection, and the retained real return
theorem, all of which hold here. The prepared-input pole realization
then supplies its ordinary pole witnesses. The independent auxiliary
returns of \cite[Theorem 3.1]{Returns} preserve these finite-local square
conditions as additional open conditions.

The source retains its nine point blocks, two pole blocks, and the same
retained roles. Applying the same finite-algebra encoding therefore
uses the same twenty-five total coordinates. This proves the theorem.
\end{proof}

\newpage
\section{Exact computational checks}

The companion checker is
\begin{center}\small
\texttt{computation/pole\_parameters/}\allowbreak
\texttt{square\_complement\_baseline\_v43/check.py}.
\end{center}
It verifies the square identity, independently enumerates all $76249$
fixed-prime residue inputs, and checks the claimed valuations on exact
integer and nonintegral examples. An in-memory substitution in the native
builder verifies the new scalar and unchanged variable roles and finite
block census; it does not modify the retained source. The certificate also
checks the loss-of-depth example above. These finite computations do not
replace the infinite arithmetic premises or establish a global rational
square-complement condition.

\begin{thebibliography}{9}
\raggedright
\bibitem{General}
\emph{A defining polynomial in twenty-five variables: combined scalar
and comparison conditions}, project source, version 19,
\texttt{papers/few\_variables25/}\allowbreak
\texttt{few\_variables25\_merged\_20261110\_v19.tex}.
\bibitem{UnitNorm}
\emph{A single norm determinant for the 25- and 13-variable constructions},
project source, version 28,
\texttt{papers/finite\_algebra/}\allowbreak
\texttt{unit\_norm\_20261114\_v28.tex}.
\bibitem{Returns}
\emph{Fixed pole fibers and independent detector indices},
project source, version 42,
\texttt{papers/pole\_parameters/}\allowbreak
\texttt{common\_parameter\_guard\_reduction\_20261116\_v42.tex}.
\bibitem{Normalization}
\emph{Square-complement normalization does not preserve the pole predicate},
project source, version 43,
\texttt{negative\_results/papers/pole\_parameters/}\allowbreak
\texttt{square\_complement\_normalization\_20261116\_v43.tex}.
\end{thebibliography}
\end{document}
